

\begingroup
\setlength{\textfloatsep}{7pt plus 2pt minus 1pt}
\setlength{\floatsep}{7pt plus 2pt minus 1pt}
\setlength{\intextsep}{6pt plus 2pt minus 1pt}
\setlength{\abovecaptionskip}{3pt}

\begin{table}[!htbp]
\centering
\begingroup
\fontsize{10.5}{11.7}\selectfont\setlength{\tabcolsep}{4.5pt}
\renewcommand{\arraystretch}{1.06}
\resizebox{\textwidth}{!}{%
\begin{tabular}{@{}lrrrrrrrrrrrrrrr@{}}
\toprule
 & \multicolumn{6}{c}{\textbf{ASR by model}} & \multicolumn{7}{c}{\textbf{ASR by attack type}} & \multicolumn{2}{c}{\textbf{TSR}} \\
\cmidrule(lr){2-7}\cmidrule(lr){8-14}\cmidrule(lr){15-16}
\textbf{Defense} & DeepSeek & Kimi & Son.~4 & Son.~4.5 & GLM-5 & \textbf{Mean} & $b(U)$ & $b(\pi)$ & $b(T)$ & $b(A)$ & $b(O)$ & $b(s)$ & \textbf{Comp.} & Benign & Attacked \\
\midrule
None & 23.2 & 17.4 & 20.6 & 11.1 & 12.9 & 17.0 & 15.0 & 21.2 & 13.7 & 19.0 & 12.4 & 18.6 & 21.8 & 95.3 & 83.7 \\
\midrule
Spotlight & 22.1 & 17.5 & 18.6 & 8.1 & 11.8 & 15.6 & 12.7 & 19.0 & 12.0 & 17.5 & 11.5 & 18.9 & 19.8 & 98.7 & 85.8 \\
Sandwich & 16.6 & 10.3 & 10.6 & 4.7 & 9.8 & 10.4 & 6.3 & 10.9 & 7.9 & 12.7 & 7.4 & 15.8 & 14.4 & 98.5 & 89.4 \\
\midrule
DRIFT-iso & 7.8 & 7.5 & 7.5 & 5.2 & 5.0 & 6.6 & 2.4 & 5.4 & 2.7 & 3.9 & 9.5 & 12.1 & 14.0 & 98.0 & 93.3 \\
\midrule
DRIFT-full & 7.2 & 6.1 & 9.0 & 6.0 & 6.6 & 7.0 & 2.5 & 5.7 & 3.9 & 4.2 & 11.4 & 11.4 & 13.3 & 88.9 & 83.4 \\
Progent & 24.8 & 14.4 & 15.6 & 9.8 & 12.6 & 15.4 & 12.4 & 18.6 & 12.0 & 17.0 & 13.2 & 18.1 & 18.4 & 74.9 & 65.0 \\
\midrule
CaMeL & 1.7 & 1.2 & 1.2 & 1.2 & 1.0 & 1.3 & 1.3 & 1.4 & 1.1 & 1.5 & 1.0 & 1.3 & 1.1 & 13.3 & 13.9 \\
\midrule
ARGUS & 14.5 & 9.6 & 10.0 & 5.7 & 9.1 & 9.8 & 8.1 & 11.1 & 8.5 & 10.2 & 7.6 & 12.3 & 11.3 & 78.9 & 69.8 \\
\bottomrule
\end{tabular}
}
\endgroup
\caption{Defense robustness and utility (\%). Model columns pool retained single-primitive and compound attacks; Mean is unweighted across models. Primitive columns average each primitive's retained single-primitive attacks across models; Comp. reports the compound attack. TSR columns are model means without an injection (benign) and under attack. Per-model results are in Appendix~\ref{app:defense-breakdown}.}
\label{tab:defense-overview}
\end{table}

\subsection{Overall defense performance}

Without defenses, GLM-5 and Sonnet 4.5 have the lowest ASRs,
12.9\% and 11.1\%, compared with
17.4--23.2\%
for the others; both generally remain among the most robust under defenses.
Among configurations retaining at least 90\% benign TSR, DRIFT-isolation has
the lowest mean ASR (6.6\%;
98.0\%/93.3\%
benign/attacked TSR), followed by Sandwich
(10.4\%;
98.5\%/89.4\%).
CaMeL reaches 1.3\% ASR but only
13.3\% benign TSR; ARGUS reaches
9.8\% ASR with 78.9\% benign TSR.

DRIFT's added plan and checklist layer does not lower ASR
(7.0\%) but lowers benign TSR
(88.9\%) relative to isolation. Progent remains
near the undefended ASR (15.4\% versus
17.0\%) while lowering benign TSR to
74.9\%. Both constrain which tool calls the agent
may make, but a tampered call uses a tool the task needs, with arguments
resolved from runtime content: they rarely block an attacked call (Section
\ref{sec:runtime}).



\subsection{Coverage beyond direct objective manipulation}

The six primitive columns of Table \ref{tab:defense-overview} disaggregate
attack success by the belief that the planted content targets while holding the
predefined harmful outcome fixed; Comp. reports the compound attack. \(b(U)\),
the objective primitive, corresponds most directly to the instruction-override
attacks common in prior evaluations.

Without defenses, \(b(\pi)\), \(b(s)\), and \(b(A)\) reach
21.2\%, 18.6\%, and 19.0\% mean ASR,
compared with 15.0\% for \(b(U)\).
This contrast is clearest for DRIFT-isolation, whose protection relies on
detecting and removing task-redirecting manipulative content: it reduces \(b(U)\) to
2.4\% ASR, but leaves
9.5\% for \(b(O)\) and
12.1\% for \(b(s)\). DRIFT-full includes the same isolator
and shows the same contrast (2.5\% versus
11.4\% and 11.4\%).
Except for CaMeL, which yields uniformly low ASR
(1.0--1.5\%),
\(b(U)\) is among the three lowest-ASR primitives under every defense.

\subsection{How agents respond to visible planted content}

The trajectory annotations (Section \ref{sec:setup}) show how GLM-5 and Sonnet
4.5 respond to planted content they can see.



\textbf{Without defenses.} On the annotated trajectories (first attack per case and primitive), GLM-5 and Sonnet 4.5 have similar ASR
(11.1\% and 10.2\%) and following
rates (12.9\% and
12.0\%) but otherwise differ:
GLM resists 25.1\% and shows no observable engagement
in 61.5\% of trajectories, whereas Sonnet resists
44.3\% and shows no observable engagement in
43.4\%.
Among resisted trajectories, both cite the task-designated authority in about
99\% of cases, but Sonnet explicitly flags manipulation more often than GLM
(43.4\% versus
18.0\%).
For both models, flagging is highest for \(b(U)\)
(37.5\% for GLM, 69.2\% for Sonnet)
and lowest for \(b(s)\) or \(b(O)\), suggesting that the models are more aware of
content targeting their objective than of content that
misstates the state or what they observe.

\textbf{Under prompt defenses.} Prompt defenses leave the planted content visible in every trajectory. They
nevertheless change the response in different ways.

\begin{figure}[!htbp]
\centering
\includegraphics[width=\textwidth]{figures/prompt_defense_behavior_grouped_by_defense.pdf}
\setlength{\abovecaptionskip}{0pt}\setlength{\parskip}{0pt}
\caption{Agent responses by primitive without a defense and under
Spotlighting and Sandwich.
Bars cover single-primitive attacks; tick letters
$U,\pi,T,A,O,s$ denote the primitives $b(U),\dots,b(s)$. Hatching marks the
resisted trajectories in which the agent explicitly identifies manipulation.}
\label{fig:prompt-defense-behavior}
\end{figure}

Both defenses lower following of \(b(U)\) and \(b(\pi)\) for both models, including
Spotlighting, whose effects on the other primitives are mixed. For GLM, following
\(b(O)\) rises from
7.1\% without a defense to
18.6\% under Spotlighting.
Sandwich more consistently lowers aggregate following (GLM:
12.9\% to 9.3\%;
Sonnet: 12.0\% to
5.4\%), but mainly shifts behavior toward
no observable engagement (61.5\% to
69.7\%;
43.4\% to
71.0\%), and flagging of manipulative content also falls.

\subsection{Runtime defenses}\label{sec:runtime}

Runtime defenses can remove content before the model reads it or reject a
resulting action. Figure \ref{fig:runtime-defense-behavior} compares what the
agent observes under DRIFT-isolation, DRIFT-full, and Progent.

\begin{figure}[!htbp]
\centering
\includegraphics[width=\textwidth]{figures/runtime_defense_impact_grouped_by_defense.pdf}
\setlength{\abovecaptionskip}{0pt}\setlength{\parskip}{0pt}
\caption{Agent-facing outcomes under runtime defenses for GLM-5 and Claude
Sonnet 4.5, grouped by defense and primitive (tick letters as in
Figure~\ref{fig:prompt-defense-behavior}). Hatching distinguishes
content stripped before action from attacked calls blocked by runtime policy.
Incomplete denotes an unfinished trajectory.}
\label{fig:runtime-defense-behavior}
\end{figure}

\begin{table}[!htbp]
\centering
\begingroup
\tablefont
\renewcommand{\arraystretch}{1.06}
\resizebox{\ifdim\width>\textwidth\textwidth\else\width\fi}{!}{%
\begin{tabular}{@{}llrrr@{}}
\toprule
\textbf{Defense} & \textbf{Model} & \begin{tabular}[b]{@{}c@{}}\textbf{Attacked} \\ \textbf{emitted}\end{tabular} & \begin{tabular}[b]{@{}c@{}}\textbf{Blocked} \\ \textbf{(rate)}\end{tabular} & \begin{tabular}[b]{@{}c@{}}\textbf{Attack} \\ \textbf{landed}\end{tabular} \\
\midrule
DRIFT-full & GLM-5 & 26 & 0 (0.0\%) & 24 \\
 & Sonnet 4.5 & 31 & 3 (9.7\%)\textsuperscript{A} & 27 \\
\midrule
Progent & GLM-5 & 64 & 5 (7.8\%)\textsuperscript{A} & 52 \\
 & Sonnet 4.5 & 52 & 3 (5.8\%)\textsuperscript{A} & 41 \\
\bottomrule
\end{tabular}
}
\endgroup
\caption{Action-level enforcement; attacked calls are counted when emitted. Superscript A marks blocks by argument-level gates (DRIFT's checklist, Progent's argument schema); no tool-level gate (DRIFT's plan-position check and alignment judge, Progent's tool allowlist) blocked an attacked call.}
\label{tab:action-enforcement}
\end{table}

\textbf{Action-level policies.} Our benchmark gives these policies limited leverage:
harmful outcomes
use tools expected for the benign task, generally by altering an argument, and
many arguments must be resolved from runtime tool results. DRIFT-full and
Progent therefore block only
0.0--9.7\%
of emitted attacked calls (Table \ref{tab:action-enforcement}).
Only argument-level gates block attacked calls.
Progent's
13.1--15.6\%
absence in Figure \ref{fig:runtime-defense-behavior}, and some under
DRIFT-full, comes not from rejecting attacked calls
but from deviation due to actions being blocked.

\textbf{Content isolation.} By contrast, DRIFT's selective protection comes
principally from its
Injection Isolator (Appendix Table \ref{tab:content-strips}), which
across both DRIFT configurations and both models fully strips
52.7--59.5\%
of examined attack carriers while redacting a legitimate result in only
1.1--4.8\%
of runs.
However, stripping varies sharply by primitive:
under isolation-only, full strips remove
78.2--85.1\%
of \(b(U)\) attacks but only
26.2--34.3\%
of \(b(O)\) and \(b(s)\) attacks.
Attacks that pass DRIFT-isolation unstripped keep their effect for \(b(O)\) and
\(b(s)\): agents follow them about as often as without a defense
(10.7\% to 14.4\% for GLM,
13.3\% to 10.7\% for Sonnet).
For the other four primitives this holds for Sonnet, but GLM follows only
3 of the
79 that pass
(3.8\%, against 14.1\% without a defense).
\(b(O)\) and \(b(s)\) thus pass the isolator more often and stay as effective when
they do.

\textbf{Does the detection gap survive wording changes?}\label{sec:detection} To test whether the gap is specific to the original wording, Figure
\ref{fig:detectability-tradeoff} compares each original single-primitive attack
with its detection-optimized rewrite and the two support-expanded conditions
(Section \ref{sec:setup}), reporting ASR without a defense, detection by the
deployed DRIFT isolator, and residual ASR. For Claude Sonnet 4.5 and GLM-5 in
all four conditions, \(b(U)\) remains among the most detectable primitives and
generally
among the least effective by ASR, especially after isolation. The pattern holds
although the rewrites were optimized against a carrier-only probe, which sees
neither the user request nor the rest of the trajectory, rather than against
the deployed isolator.

\begin{figure}[!htbp]
\centering
\includegraphics[width=\textwidth]{figures/detectability_three_metric_sonnet_glm.pdf}
\setlength{\abovecaptionskip}{0pt}\setlength{\parskip}{0pt}
\caption{ASR without a defense, runtime detection, and residual ASR under
DRIFT-isolation across the four conditions. Detection counts partial or
complete removal. Gray lines span a primitive's four conditions; shading marks
\(b(O)\) and \(b(s)\). Appendix \ref{app:detection} gives the other two models.}
\label{fig:detectability-tradeoff}
\end{figure}

\endgroup
